\documentclass[aspectratio=169]{beamer}
\usetheme{Madrid}
\usepackage[utf8]{inputenc}
\usepackage[T1]{fontenc}
\usepackage{tikz}
\usepackage{amssymb}
\usepackage{graphicx}
\setbeamertemplate{navigation symbols}{}
\setbeamercolor{title}{fg=white,bg=blue!65!cyan}
\setbeamercolor{frametitle}{fg=white,bg=blue!65!cyan}
\setbeamercolor{block title}{bg=blue!65!cyan!20,fg=black}
\setbeamercolor{block body}{bg=blue!65!cyan!8,fg=black}
\setbeamercolor{item}{fg=blue!65!cyan}
\setbeamercolor{itemize item}{fg=blue!65!cyan}
\setbeamertemplate{itemize item}[triangle]

\title{For Loops}
\subtitle{range() — and the "second number is exclusive" trap}
\author{ByteForge --- Grade 8 --- App Builder}
\date{}

\begin{document}

\begin{frame}
    \titlepage
\end{frame}

\begin{frame}{Today's Plan}
    \begin{itemize}
        \item Count to Five
        \item[$\bigstar$] BOSS --- Count 1 to 5, For Real
    \end{itemize}
\end{frame}

\begin{frame}{\texttt{range(5)} stops BEFORE 5}
    \large
    \begin{itemize}
        \item \texttt{range(5)} produces 0, 1, 2, 3, 4
        \item Five numbers total, but the number 5 itself is never reached. This exact trap catches nearly every student the first time.
    \end{itemize}
\end{frame}


\begin{frame}{\texttt{range(start, stop)} — stop is still exclusive}
    \large
    \begin{itemize}
        \item \texttt{range(1, 6)} starts at 1 and stops before 6
        \item Landing exactly on 5 as the last number actually printed. To include N, the stop argument must be N+1.
    \end{itemize}
\end{frame}


\begin{frame}{BOSS — applying the same trap deliberately}
    \large
    \begin{itemize}
        \item Printing 1 through 5 inclusive requires spotting that the stop number needs to be one MORE than the last number you actually want.
    \end{itemize}
\end{frame}

\begin{frame}{Real-World Examples}
    \large
    \begin{itemize}
        \item This exact "off-by-one" mistake (using \texttt{range(1,10)} when you meant to include 10) is a genuinely common real bug that professional developers still make
        \item Not just a beginner trap.
    \end{itemize}
\end{frame}


\begin{frame}{Let Them Discover}
    \Large
    Before running anything: does \texttt{range(1, 11)} include the number 11? Write your answer down first, then check it — most people guess wrong the first time.
\end{frame}

\begin{frame}{Your Turn}
    \Large
    Open the student map, find \textbf{For Loops}, and work through it
    hands-on --- everything you just saw, they now do themselves.
\end{frame}

\end{document}
